\documentclass[11pt]{amsart}
\usepackage[utf8]{inputenc}
\usepackage[margin=1in, headheight=14.5pt]{geometry}
\usepackage{math-hw}
\pagestyle{fancy}
\fancyhead[L]{Math 918}
%\fancyhead[C]{MATH 665}
\fancyhead[R]{HW 4}

\usepackage[margin=1in, headheight=14.5pt]{geometry}
\usepackage{parskip}

%%%%%%%%%%%%%%%%%%%%%%%%%%%%%%%%%%
% fonts
\usepackage{calligra}
\usepackage{mathrsfs}
\usepackage[mathscr]{euscript}
\usepackage{bbm}

\usepackage{enumitem}
\usepackage{hyperref}
%\usepackage[colorlinks,  citecolor=green, linkcolor = black, urlcolor = cyan]{hyperref} % hyperlinks for cross-referencing; colorlinks is recommended because otherwise the boxes obscure the roman numerals; linkcolor is changed to blue after the table of contents (see firstpages.tex). As of 2021, red is not allowed by Rackham.
\usepackage[nameinlink]{cleveref}
%\usepackage{backref}
\usepackage{url} % allow hyperlink in reference
\usepackage[style=alphabetic,url=false,backref=true]{biblatex}
\addbibresource{zotero-export.bib}
\setlength{\droptitle}{-4em}


\newtoggle{solution}
\toggletrue{solution}
%\togglefalse{solution}


\iftoggle{solution}{
\includeversion{solnenv}
\includeversion{commenv}
}{
\excludeversion{solnenv}
\excludeversion{commenv}
}





\title{Homework 4}
%\author{Anna Brosowsky}
\date{\vspace{-4em}}

\begin{document}


\maketitle\thispagestyle{fancy}

Remember you are allowed to discuss with classmates (or an AI tool), but that you need to tell me who/what you discussed with + the final submitted writeup should be your own work.


\begin{enumerate}[itemsep=3em]

\item (Exercise 4.14 of [SS]) Write $K(R)$ for the total ring of fractions. A reduced Noetherian ring $R$ of characteristic $p > 0$ is called
\emph{weakly normal} if $x\in K(R)$ and $x^p \in R$ implies that $x\in R$. Any ring is
called \emph{seminormal} if $x\in K(R)$, and $x^2, x^3 \in R$ implies that $x \in R$. 
\begin{enumerate}
\item Show
that Frobenius split rings are weakly normal, WITHOUT using the fact that F-injective rings are weakly normal.\footnote{...because that fact makes this problem boring.}



\item Show that weakly normal rings
are seminormal. 

\end{enumerate}


\item (Exercise 4.12 of [SS]) Let $R$ be a Noetherian char $p$ ring. Suppose that $\varphi:F_*^eR\to R$ is an $R$-linear map. 
\begin{enumerate}
\item Fix $I$ an ideal. Prove that $\varphi(F_*^e(\Gamma_I(R)))\subset \Gamma_I(R)$, where $\Gamma_I(R) = \bigcup_{t\in \N} \ann(I^t)$.
\item
Suppose further that $R$ is reduced. Prove that if $Q \subset R$ is a minimal prime of $R$, then $\varphi(F_*^eQ)\subset Q$.
\end{enumerate}



\item (Exercise 23 of [MP]) Let $(R,\mf m,k)$ be a local Noetherian ring of char $p$ and $d=\dim R$. 

\begin{enumerate}
\item Prove that $H^d_{\mf m}(R)\ne 0$. \footnote{Hint: You may find it helpful to use the Theorem from Tuesday 03/10, as well as to use facts about completion from 906 (especially facts from Section~3.4 of this semester's notes!)}

\item Prove that the kernel of the natural Frobenius action on $H^d_{\mf m}(R)$ is a proper submodule of $H^d_{\mf m}(R)$.
\end{enumerate}


\item (Exercise 7.15 of [SS]) Let $(R,\mf m, k)$ be a local Cohen-Macaulay ring\footnote{Fact (you may use without proof; possibly helpful for (c) implies (a)): If $R$ is a CM ring and $\underline x$ and $\underline y$ are sop's, then $\dim_k(\frac{\langle x_1,\ldots, x_d\rangle:\mf m}{\langle x_1,\ldots, x_d\rangle}) =\dim_k(\frac{\langle y_1,\ldots, y_d\rangle:\mf m}{\langle y_1,\ldots, y_d\rangle})$. This number is the \emph{Cohen-Macaulay type} of $R$.} of dimension $d$. Prove that TFAE:
\begin{enumerate}
\item The natural Frobenius action on $H^d_{\mf m}(R)$ is injective (in other words, $R$ is $F$-injective)
\item For every s.o.p. $f_1,\ldots, f_d$, the ideal $I=\langle f_1,\ldots, f_d\rangle $ has
\[
I = \langle g \: | \: g^p \in I^{[p]}\rangle.
\]
\item There exists an s.o.p. $f_1,\ldots, f_d$, such that the ideal $I=\langle f_1,\ldots, f_d\rangle $ has\footnote{Hint: Prove for any zero-dimensional local ring $(S,\mf m)$ that the extension $(0:\mf m)\hookrightarrow S$ is essential.}
\[
I = \langle g \: | \: g^p \in I^{[p]}\rangle.
\]
\end{enumerate}

\end{enumerate}

\end{document}
